% GENERATED FILE. Edit canonical lessons, then run npm run generate:books.
% canonical-source-hash: fnv1a64:f03b6a4593f7cf2d
% canonical-lessons: SA-C464-vacalah, SA-C464-mitabhasi, SA-C464-sambodhayati, SA-C464-akrosati, SA-C464-pralapah

\chapter{Talkers and Quiet Ones}
\label{ch:sa-talkers-and-quiet-ones}
\hlchaptermodality{464}

\begin{chapteropening}
\textbf{By the end of this chapter:} \emph{I can say whether someone talks a lot or keeps quiet.}

Complete the last lesson of chapter 464: I can say whether someone talks a lot or keeps quiet.
\end{chapteropening}

\section[\texorpdfstring{\sk{वाचालः} (vācālaḥ)}{वाचालः (vācālaḥ)}]{\sk{वाचालः} (vācālaḥ) — talkative}

\label{lesson:SA-C464-vacalah}



\begin{quote}
Before the new one: say the Sanskrit for trembles with fear, then the Sanskrit for bold.
\end{quote}

\subsection*{You'll want to know: \sk{वाचालः}}
\textbf{\sk{वाचालः}} — \emph{vācālaḥ} — \textquotedblleft{}talkative\textquotedblright{}.

\begin{quote}
\sk{वाचालः} \sk{नापितः} \sk{सदा} \sk{जल्पति।}
\end{quote}

\begin{quote}
\emph{vācālaḥ nāpitaḥ sadā jalpati.}
\end{quote}

\begin{quote}
The talkative barber is always chattering.
\end{quote}

\begin{cousinweb}
\textbf{\sk{वाचालः}} (\emph{vācālaḥ}) is built on \emph{vāc}, \textquotedblleft{}speech, voice\textquotedblright{}; MW gives \textquotedblleft{}talkative, chattering\textquotedblright{} and \textquotedblleft{}boasting\textquotedblright{} (MW).
\end{cousinweb}

\begin{grammarlens}[title={What you've built}]
One more word for describing what you lived and felt.
\end{grammarlens}

\subsection*{Your turn}
\begin{itemize}
\raggedright
  \item \emph{Say it:} \emph{vācālaḥ}
  \item \emph{Say it:} \emph{vācālaḥ}, once more
  \item \emph{Say it:} \emph{vācālaḥ nāpitaḥ sadā jalpati.} — The talkative barber is always chattering.
  \item \emph{Recall:} say the Sanskrit for trembles with fear, then the Sanskrit for bold, then say \emph{vācālaḥ} again
\end{itemize}

\begin{tcolorbox}[breakable,colback=teal!4,colframe=teal!35!black,title={Before you move on}]
What does \sk{वाचालः} mean? (\textquotedblleft{}Talkative\textquotedblright{}.) Say it once more.
\end{tcolorbox}
\section[\texorpdfstring{\sk{मितभाषी} (mitabhāṣī)}{मितभाषी (mitabhāṣī)}]{\sk{मितभाषी} (mitabhāṣī) — of few words}

\label{lesson:SA-C464-mitabhasi}



\begin{quote}
Before the new one: say the Sanskrit for bold, then the Sanskrit for talkative.
\end{quote}

\subsection*{You'll want to know: \sk{मितभाषी}}
\textbf{\sk{मितभाषी}} — \emph{mitabhāṣī} — \textquotedblleft{}of few words\textquotedblright{}.

\begin{quote}
\sk{मम} \sk{पिता} \sk{मितभाषी} \sk{अस्ति।}
\end{quote}

\begin{quote}
\emph{mama pitā mitabhāṣī asti.}
\end{quote}

\begin{quote}
My father is a man of few words.
\end{quote}

\begin{cousinweb}
\textbf{\sk{मितभाषी}} (\emph{mitabhāṣī}) is built from \emph{mita}, \textquotedblleft{}measured\textquotedblright{}, and \emph{bhāṣin}, \textquotedblleft{}speaking\textquotedblright{}, from the root √bhāṣ of \textbf{\sk{भाषते}}: one who speaks by measure (MW).
\end{cousinweb}

\begin{grammarlens}[title={What you've built}]
One more word for describing what you lived and felt.
\end{grammarlens}

\subsection*{Your turn}
\begin{itemize}
\raggedright
  \item \emph{Say it:} \emph{mitabhāṣī}
  \item \emph{Say it:} \emph{mitabhāṣī}, once more
  \item \emph{Say it:} \emph{mama pitā mitabhāṣī asti.} — My father is a man of few words.
  \item \emph{Recall:} say the Sanskrit for bold, then the Sanskrit for talkative, then say \emph{mitabhāṣī} again
\end{itemize}

\begin{tcolorbox}[breakable,colback=teal!4,colframe=teal!35!black,title={Before you move on}]
What does \sk{मितभाषी} mean? (\textquotedblleft{}Of few words\textquotedblright{}.) Say it once more.
\end{tcolorbox}
\section[\texorpdfstring{\sk{सम्बोधयति} (sambodhayati)}{सम्बोधयति (sambodhayati)}]{\sk{सम्बोधयति} (sambodhayati) — calls out to; addresses}

\label{lesson:SA-C464-sambodhayati}



\begin{quote}
Before the new one: say the Sanskrit for talkative, then the Sanskrit for of few words.
\end{quote}

\subsection*{You'll want to know: \sk{सम्बोधयति}}
\textbf{\sk{सम्बोधयति}} — \emph{sambodhayati} — \textquotedblleft{}calls out to; addresses\textquotedblright{}.

\begin{quote}
\sk{माता} \sk{द्वारात्} \sk{पुत्रं} \sk{सम्बोधयति।}
\end{quote}

\begin{quote}
\emph{mātā dvārāt putraṃ sambodhayati.}
\end{quote}

\begin{quote}
Mother calls out to her son from the door.
\end{quote}

Past (she …): \sk{सा} \sk{पुत्रं} \sk{सम्बोधितवती}, \emph{sā putraṃ sambodhitavatī}.

\begin{cousinweb}
\textbf{\sk{सम्बोधयति}} (\emph{sambodhayati}) is built on \emph{sam} and the causative of the root √budh, \textquotedblleft{}to wake, to know\textquotedblright{}, the root of \textbf{\sk{बुध्यते}}: to rouse someone's notice; MW gives the causative the senses \textquotedblleft{}to cause to wake up\textquotedblright{}, \textquotedblleft{}to inform, advise\textquotedblright{} and \textquotedblleft{}to call to\textquotedblright{} (MW).
\end{cousinweb}

\begin{grammarlens}[title={What you've built}]
One more word for describing what you lived and felt.
\end{grammarlens}

\subsection*{Your turn}
\begin{itemize}
\raggedright
  \item \emph{Say it:} \emph{sambodhayati}
  \item \emph{Say it:} \emph{sambodhayati}, once more
  \item \emph{Say it:} \emph{mātā dvārāt putraṃ sambodhayati.} — Mother calls out to her son from the door.
  \item \emph{Recall:} say the Sanskrit for talkative, then the Sanskrit for of few words, then say \emph{sambodhayati} again
\end{itemize}

\begin{tcolorbox}[breakable,colback=teal!4,colframe=teal!35!black,title={Before you move on}]
What does \sk{सम्बोधयति} mean? (\textquotedblleft{}Calls out to; addresses\textquotedblright{}.) Say it once more.
\end{tcolorbox}
\section[\texorpdfstring{\sk{आक्रोशति} (ākrośati)}{आक्रोशति (ākrośati)}]{\sk{आक्रोशति} (ākrośati) — reviles; scolds loudly}

\label{lesson:SA-C464-akrosati}



\begin{quote}
Before the new one: say the Sanskrit for of few words, then the Sanskrit for calls out to.
\end{quote}

\subsection*{You'll want to know: \sk{आक्रोशति}}
\textbf{\sk{आक्रोशति}} — \emph{ākrośati} — \textquotedblleft{}reviles; scolds loudly\textquotedblright{}.

\begin{quote}
\sk{कुपितः} \sk{स्वामी} \sk{सेवकम्} \sk{आक्रोशति।}
\end{quote}

\begin{quote}
\emph{kupitaḥ svāmī sevakam ākrośati.}
\end{quote}

\begin{quote}
The angry master reviles his servant.
\end{quote}

Past (he …): \sk{सः} \sk{सेवकम्} \sk{आक्रुष्टवान्}, \emph{saḥ sevakam ākruṣṭavān}.

\begin{cousinweb}
\textbf{\sk{आक्रोशति}} (\emph{ākrośati}) is built from \emph{ā} and the root √kruś, \textquotedblleft{}to cry out\textquotedblright{}; MW gives \textquotedblleft{}to call to anyone in an abusive manner, scold at, revile\textquotedblright{} (MW).
\end{cousinweb}

\begin{grammarlens}[title={What you've built}]
One more word for describing what you lived and felt.
\end{grammarlens}

\subsection*{Your turn}
\begin{itemize}
\raggedright
  \item \emph{Say it:} \emph{ākrośati}
  \item \emph{Say it:} \emph{ākrośati}, once more
  \item \emph{Say it:} \emph{kupitaḥ svāmī sevakam ākrośati.} — The angry master reviles his servant.
  \item \emph{Recall:} say the Sanskrit for of few words, then the Sanskrit for calls out to, then say \emph{ākrośati} again
\end{itemize}

\begin{tcolorbox}[breakable,colback=teal!4,colframe=teal!35!black,title={Before you move on}]
What does \sk{आक्रोशति} mean? (\textquotedblleft{}Reviles; scolds loudly\textquotedblright{}.) Say it once more.
\end{tcolorbox}
\section[\texorpdfstring{\sk{प्रलापः} (pralāpaḥ)}{प्रलापः (pralāpaḥ)}]{\sk{प्रलापः} (pralāpaḥ) — prattle}

\label{lesson:SA-C464-pralapah}



\begin{quote}
Before the new one: say the Sanskrit for calls out to, then the Sanskrit for reviles.
\end{quote}

\subsection*{You'll want to know: \sk{प्रलापः}}
\textbf{\sk{प्रलापः}} — \emph{pralāpaḥ} — \textquotedblleft{}prattle\textquotedblright{}.

\begin{quote}
\sk{तस्य} \sk{प्रलापं} \sk{कः} \sk{शृणोति}?
\end{quote}

\begin{quote}
\emph{tasya pralāpaṃ kaḥ śṛṇoti?}
\end{quote}

\begin{quote}
Who listens to his prattle?
\end{quote}

A masculine noun.

\begin{cousinweb}
\textbf{\sk{प्रलापः}} (\emph{pralāpaḥ}) is built from \emph{pra} and the root √lap, \textquotedblleft{}to talk\textquotedblright{}; MW gives \textquotedblleft{}talk, prattling, chattering\textquotedblright{} and also \textquotedblleft{}lamentation\textquotedblright{} (MW).
\end{cousinweb}

\begin{grammarlens}[title={What you've built}]
One more word for describing what you lived and felt.
\end{grammarlens}

\subsection*{Your turn}
\begin{itemize}
\raggedright
  \item \emph{Say it:} \emph{pralāpaḥ}
  \item \emph{Say it:} \emph{pralāpaḥ}, once more
  \item \emph{Say it:} \emph{tasya pralāpaṃ kaḥ śṛṇoti?} — Who listens to his prattle?
  \item \emph{Recall:} say the Sanskrit for calls out to, then the Sanskrit for reviles, then say \emph{pralāpaḥ} again
\end{itemize}

\begin{tcolorbox}[breakable,colback=teal!4,colframe=teal!35!black,title={Before you move on}]
What does \sk{प्रलापः} mean? (\textquotedblleft{}Prattle\textquotedblright{}.) Say it once more.
\end{tcolorbox}
